\setcounter{aufgabe}{40}
\hypertarget{e5}{}\einheitenkopf{Einheit 5 von 5 · Anwendungen}
\zweigzeile{Hier lernst du, Tarife und andere Sachsituationen mit linearen Funktionen zu beschreiben, zu berechnen und zu vergleichen · neu in diesem Jahr · P10 oft · baut auf: Gleichung, Graph und Schnittpunkt (Einheit 2 bis 4), Dreisatz}

\begin{aufgabe}{\textbf{Ich finde Anfangswert und Änderung im Text.} Unterstreiche den Anfangswert. Kreise die Änderung je Einheit ein.}
\begin{teile}
\teil Ein Handyvertrag kostet 8 € Grundgebühr im Monat und 9 Cent je Gesprächsminute.
\teil In einem Wassertank sind 500 Liter. Jede Stunde werden 20 Liter entnommen.
\teil Ein Fitnessstudio verlangt einmalig 29 € Aufnahmegebühr und dann 15 € pro Monat.
\teil Eine Taxifahrt kostet 2,10 € je Kilometer. Dazu kommt ein Grundpreis von 3,90 €.
\teil Ein Baby ist bei der Geburt 50 cm groß und wächst in den ersten Monaten etwa 3 cm pro Monat.
\end{teile}
\end{aufgabe}

\begin{aufgabe}{\textbf{Ich kann die passende Gleichung zu einem Tarif erkennen.} Kreuze die Gleichung an, die den Preis $y$ in € beschreibt.}
\begin{teile}
\teil Carsharing: 6 € Grundgebühr und 0,30 € je Kilometer ($x$: Strecke in km). \\
\kreuz{$y = 6x + 0{,}3$} \\ \kreuz{$y = 0{,}3x + 6$} \\ \kreuz{$y = 6{,}3x$}
\teil Kletterpark: 12 € Eintritt und 4 € je Stunde für die Leihausrüstung ($x$: Zeit in Stunden). \\
\kreuz{$y = 12x + 4$} \\ \kreuz{$y = 16x$} \\ \kreuz{$y = 4x + 12$}
\teil Strom: Grundpreis 9 € im Monat und 30 Cent je Kilowattstunde ($x$: verbrauchte Kilowattstunden). (P10 2021) \\
\kreuz{$y = 30x + 9$} \\ \kreuz{$y = 0{,}3x + 9$} \\ \kreuz{$y = 9x + 30$}
\end{teile}
\end{aufgabe}

\begin{aufgabe}{\textbf{Ich kann eine Gleichung zu einer Sachsituation aufstellen.} Stelle die Gleichung auf.}
\begin{teile}
\teil Ein Fahrradverleih verlangt 5 € Grundgebühr und 3 € je Stunde ($x$: Zeit in Stunden, $y$: Preis in €). \\ \feldl{y}
\teil Ein Tank mit 800 Litern Heizöl verliert jeden Tag 25 Liter ($x$: Zeit in Tagen, $y$: Heizöl in Litern). \\ \feldl{y}
\teil Ein Kopierladen verlangt 8 Cent je Kopie und einmalig 2 € für die Kopierkarte ($x$: Anzahl der Kopien, $y$: Preis in €). \\ \feldl{y}
\end{teile}
\end{aufgabe}

\begin{aufgabe}{\textbf{Ich kann berechnen, welcher Wert nach einer bestimmten Zeit erreicht ist.} Berechne.}
\begin{teile}
\teil In einem Sparschwein sind 40 €. Jede Woche kommen 6 € dazu. Wie viel Geld ist nach 15 Wochen darin?
\teil Ein Schwimmbecken enthält 60 m³ Wasser. Beim Ablassen fließen jede Stunde 7,5 m³ ab. Wie viel Wasser ist nach 6 Stunden noch im Becken?
\teil Der Wasserstand eines Sees sinkt im Sommer gleichmäßig. Wie hoch ist er nach 10 Wochen? \par
\wertetabelle[120,116,112,108]{\text{Woche}}{\text{Stand in cm}}{0,1,2,3}
\teil Lara hat 350 € gespart und legt jeden Monat 45 € dazu. Nach wie vielen Monaten hat sie zum ersten Mal mehr als 1000 €? (P10 2022)
\end{teile}
\end{aufgabe}

\begin{aufgabe}{\textbf{Ich kann zwei Tarife für eine bestimmte Nutzung vergleichen.} Berechne die Kosten und entscheide.}
\begin{teile}
\teil Handytarif A: 6 € Grundgebühr im Monat und 8 Cent je Minute. Handytarif B: keine Grundgebühr, 15 Cent je Minute. Welcher Tarif ist bei 100 Minuten im Monat günstiger?
\teil Tarif C: 12 € im Monat, 50 Minuten sind frei, jede weitere Minute kostet 20 Cent. Tarif D: 4 € im Monat und 10 Cent je Minute. Welcher Tarif ist bei 120 Minuten im Monat günstiger? Begründe mit einer Rechnung. (P10 2023)
\end{teile}
\end{aufgabe}

\begin{aufgabe}{\textbf{Ich kann berechnen, ab wann ein Tarif günstiger ist.} Ein Bootsverleih A verlangt 8 € Grundgebühr und 2,50 € je Stunde. Bootsverleih B verlangt 4,50 € je Stunde ohne Grundgebühr. $x$ ist die Zeit in Stunden, $y$ der Preis in €.}
\begin{teile}
\teil Stelle für beide Verleihe die Gleichung auf. \quad A: \feldl{y} \quad B: \feldl{y}
\teil Ab welcher Mietdauer ist Verleih A günstiger? Berechne. (P10 2016)
\teil Bleibt A auch bei 10 Stunden günstiger? Begründe.
\end{teile}
\end{aufgabe}

\begin{aufgabe}{\textbf{Ich finde den Fehler beim Aufstellen einer Tarifgleichung.} Eine Musikschule verlangt 20 € Anmeldegebühr und 35 € je Monat. Tim stellt die Gleichung auf und berechnet die Kosten für 6 Monate:}
\rechnung{y &= 20x + 35 \\ y &= 20 \cdot 6 + 35 = 155}
\begin{teile}
\teil Finde den Fehler, benenne ihn und korrigiere die Rechnung.
\teil Ein Paketdienst verlangt 3,50 € Grundpreis und 60 Cent je Kilogramm. Stelle die Gleichung auf und berechne den Preis für ein Paket mit 8 kg.
\end{teile}
\end{aufgabe}

\begin{aufgabe}{\textbf{Ich kann einem Tarif seinen Graphen zuordnen.} Drei Anbieter verleihen E-Scooter. Anbieter A: 2 € Entsperrgebühr und 0,20 € je Minute. Anbieter B: 0,30 € je Minute, keine Entsperrgebühr. Anbieter C: 5 € Grundgebühr und 0,10 € je Minute.}

\begin{ksys}[xmin=0,xmax=40,ymin=0,ymax=12,xstep=5,ystep=1,xlabel=$x$ in min,ylabel=$y$ in €,ablesen]
\gerade[36]{0.3}{0}{\mathrm{I}}\gerade[12]{0.1}{5}{\mathrm{II}}\gerade[6]{0.2}{2}{\mathrm{III}}
\end{ksys}

\begin{teile}
\teil Ordne jedem Anbieter seinen Graphen zu. \\ A: \leerfeld \quad B: \leerfeld \quad C: \leerfeld
\teil Begründe deine Zuordnung für Anbieter B. (P10 2023)
\teil Lies ab: Welcher Anbieter ist bei einer Fahrt von 10 Minuten am günstigsten?
\end{teile}
\end{aufgabe}

\begin{aufgabe}{\textbf{Ich kann zu einer Gleichung eine passende Situation beschreiben.} Beschreibe eine Situation, zu der die Gleichung passt. Sage, wofür $x$ und $y$ stehen und was die beiden Zahlen bedeuten.}
\begin{teile}
\teil $y = 2{,}5x + 10$
\teil $y = -15x + 300$
\teil $y = 0{,}12x$
\end{teile}
\end{aufgabe}
